% Part: sets-functions-relations
% Chapter: relations
% Section: orders

\documentclass[../../../include/open-logic-section]{subfiles}

\begin{document}

\olfileid{sfr}{rel}{ord}
\olsection{ترتیباں (Orders)}

\begin{explain}
ساڈے کئی موازنے کجھ چیزاں نوں، کسے خاص لحاظ نال، ہور چیزاں توں
«گھٹ»، اوہناں دے «برابر»، یا اوہناں توں «ودھ» دسن نال تعلق رکھدے نیں۔
ایہناں وچ \emph{ترتیبی تعلق (order relations)} آؤندے نیں۔ پر ترتیبی تعلقاں دیاں وکھ وکھ
قسماں نیں۔ مثال دے طور تے، کجھ وچ ضروری اے کہ کوئی وی دو چیزاں آپس وچ
قابلِ موازنہ ہون، پر کجھ وچ نہیں۔ کجھ اکو ہون (identity) نوں شامل کردے نیں
(جیویں~$\le$) تے کجھ اوہنوں کڈھ دیندے نیں (جیویں~$<$)۔
ایتھے ایہناں دی درجہ بندی (taxonomy) ساڈے کم آوے گی۔
\end{explain}

\begin{defn}[پیش ترتیب (Preorder)]
اک تعلق نوں، جے اوہ انعکاسی (reflexive) تے متعدی (transitive) دونویں ہووے،
\emph{پیش ترتیب (preorder)} آکھیا جاندا اے۔
\end{defn}

\begin{defn}[جزوی ترتیب (Partial order)]
اک پیش ترتیب جیہڑی ضد تناظری (anti-symmetric) وی ہووے، \emph{جزوی ترتیب
(partial order)} آکھلاندی اے۔
\end{defn}

\begin{defn}[خطی ترتیب (Linear order)]\ollabel{def:linearorder}
اک جزوی ترتیب جیہڑی باہم قابلِ موازنہ (connected) وی ہووے، \emph{کُل ترتیب
(total order)} یا \emph{خطی ترتیب (linear order)} آکھلاندی اے۔
\end{defn}

\begin{ex}
ہر خطی ترتیب جزوی ترتیب وی ہوندی اے، تے ہر جزوی ترتیب پیش ترتیب وی
ہوندی اے، پر اُلٹ گلاں درست نہیں۔ $A$ اُتے ہمہ گیر تعلق (universal relation) پیش ترتیب اے، کیوں
جے اوہ انعکاسی تے متعدی اے۔ پر جے $A$ وچ اک توں ودھ !!{element} ہووے،
تاں ہمہ گیر تعلق ضد تناظری نہیں ہوندا، تے ایس لئی
جزوی ترتیب وی نہیں۔
\end{ex}

\begin{ex}
\emph{لمبائی وچ ودھ نہ ہون (no longer than)} دے تعلق $\preccurlyeq$ اُتے غور
کرو، جیہڑا $\Bin^*$ اُتے اے: $x \preccurlyeq y$ اُدوں تے صرف اُدوں جدوں
$\len{x} \le \len{y}$۔
ایہ پیش ترتیب اے (انعکاسی تے متعدی)، بلکہ باہم قابلِ موازنہ وی اے، پر
جزوی ترتیب نہیں، کیوں جے ایہ ضد تناظری نہیں۔ مثال لئی، $01
\preccurlyeq 10$ تے $10 \preccurlyeq 01$، پر $01 \neq 10$۔
\end{ex}

\begin{ex}
اک اہم جزوی ترتیب، سیٹاں دے کسے سیٹ اُتے تعلق $\subseteq$ اے۔ ایہ عام
طور تے خطی ترتیب نہیں، کیوں جے جے $a \neq b$ ہووے تے اسی
$\Pow{\{a, b\}} = \{\emptyset, \{a\}, \{b\}, \{a,b\}\}$ اُتے غور کریے،
تاں ویکھدے آں کہ $\{a\} \nsubseteq \{b\}$ تے $\{a\} \neq \{b\}$ تے
$\{b\} \nsubseteq \{a\}$۔
\end{ex}

\begin{ex}
\emph{بغیر باقی پورا ونڈن (divisibility without remainder)} دا تعلق سانوں اک ایہو جیہی جزوی ترتیب
دیندا اے جیہڑی خطی ترتیب نہیں۔ صحیح عدداں $n$ تے~$m$ لئی، اسی
$n \mid m$ ایہ مطلب دسن لئی لکھدے آں کہ $n$، عدد $m$ نوں پورا ونڈدا
اے، یعنی اُدوں تے صرف اُدوں جدوں کوئی صحیح عدد~$k$ ایہو جیہا ہووے کہ
$m = kn$۔ $\Nat$ اُتے ایہ جزوی ترتیب اے، پر خطی ترتیب نہیں: مثال لئی،
$2 \nmid 3$ تے نالے $3 \nmid 2$۔ جے $\Int$ اُتے تعلق دے طور تے ویکھیے،
تاں پورا ونڈن دا تعلق صرف پیش ترتیب اے، کیوں جے ایہ ضد تناظری نہیں:
$1 \mid -1$ تے $-1 \mid 1$، پر $1 \neq -1$۔
\end{ex}

\begin{ex}
سلسلیاں دے سیٹ~$A^*$ اُتے \emph{ودھاؤ (extension)} دا تعلق ایہ اے: $s \sqsubseteq s'$
اُدوں تے صرف اُدوں جدوں $s = \emptyseq$ (خالی سلسلہ)، $s = s'$، یا
$s = \tuple{s_1, \dots, s_n}$ تے $s' =
\tuple{s_1, \dots, s_n, s_{n+1}, \dots, s_m}$ ہووے۔ جے $s \sqsubseteq s'$
ہووے، تاں اسی ایہ وی آکھدے آں کہ $s$، سلسلے~$s'$ دا \emph{مُڈھلا ٹکڑا
(initial segment)}
اے۔ $A^*$ اُتے ودھاؤ دا تعلق جزوی ترتیب اے پر خطی ترتیب نہیں؛ مثال لئی،
جے $a \neq b$ ہووے، تاں $ab \not\sqsubseteq ba$ تے $ba \not\sqsubseteq ab$۔
\end{ex}

\begin{defn}[سخت ترتیب (Strict order)]
\emph{سخت ترتیب (strict order)} اک اوہ تعلق اے جیہڑا اپنے آپ نال تعلق
نہ ہون والا (irreflexive)، یک طرفہ (asymmetric)، تے متعدی (transitive) ہوندا اے۔
\end{defn}

\begin{defn}[سخت خطی ترتیب (Strict linear order)]\ollabel{def:strictlinearorder}
اک سخت ترتیب جیہڑی باہم قابلِ موازنہ وی ہووے، \emph{سخت کُل ترتیب
(strict total order)} یا \emph{سخت خطی ترتیب (strict linear order)} آکھلاندی اے۔
\end{defn}

\begin{ex}
$\le$ اوہ خطی ترتیب اے جیہڑی سخت خطی ترتیب~$<$ دے مطابق اے۔
$\subseteq$ اوہ جزوی ترتیب اے جیہڑی سخت ترتیب~$\subsetneq$ دے مطابق اے۔
\end{ex}

کوئی وی سخت ترتیب $R$ جیہڑی~$A$ اُتے ہووے، اوہدے وچ قطر (diagonal) $\Id{A}$ شامل
کر کے، یعنی سارے جوڑے~$\tuple{x, x}$ شامل کر کے، جزوی ترتیب وچ بدلی
جا سکدی اے۔ (ایہنوں $R$ دی \emph{انعکاسی تکمیل (reflexive closure)} آکھدے نیں۔) اُلٹ پاسے،
جزوی ترتیب توں شروع کر کے، $\Id{A}$ کڈھ کے سخت ترتیب حاصل کیتی جا
سکدی اے۔ اگلے دو نتیجے ایہ گل ٹھیک ٹھیک بیان کردے نیں۔

\begin{prop}\ollabel{prop:stricttopartial}
جے $R$، سیٹ~$A$ اُتے سخت ترتیب اے، تاں $R^+ = R \cup \Id{A}$ جزوی
ترتیب اے۔ نالے، جے $R$ سخت خطی ترتیب اے، تاں $R^+$ خطی ترتیب اے۔
\end{prop}

\begin{proof}
فرض کرو کہ $R$ سخت ترتیب اے، یعنی $R \subseteq A^2$ تے $R$ اپنے آپ نال
تعلق نہ ہون والا، یک طرفہ، تے متعدی اے۔ رکھو کہ $R^+ = R \cup \Id{A}$۔
سانوں ثابت کرنا اے کہ $R^+$ انعکاسی، ضد تناظری، تے متعدی اے۔

$R^+$ صاف طور تے انعکاسی اے، کیوں جے $\tuple{x, x} \in \Id{A} \subseteq
R^+$ ہر $x \in A$ لئی درست اے۔

ایہ ثابت کرن لئی کہ $R^+$ ضد تناظری اے، برہانِ خلف (reductio) لئی فرض کرو کہ $R^+xy$
تے $R^+yx$، پر $x \neq y$۔ کیوں جے $\tuple{x,y} \in R \cup \Id{A}$، پر
$\tuple{x, y} \notin \Id{A}$، ایس لئی لازماً $\tuple{x, y} \in R$، یعنی
$Rxy$۔ ایس طرح~$Ryx$ وی۔ پر ایہ ایس مفروضے دے اُلٹ اے کہ $R$ یک طرفہ اے۔

تعدیت ثابت کرن لئی، من لو کہ $R^+xy$ تے $R^+yz$۔ جے $\tuple{x, y} \in R$
تے $\tuple{y,z} \in R$ دونویں ہون، تاں $\tuple{x, z} \in R$، کیوں جے
$R$~متعدی اے۔ نہیں تاں، یا $\tuple{x, y} \in \Id{A}$، یعنی $x = y$، یا
$\tuple{y, z} \in \Id{A}$، یعنی $y = z$۔ پہلی صورت وچ، مفروضے مطابق
$R^+yz$، تے $x = y$، سو $R^+xz$۔ دوسری صورت وچ وی ایس طرح۔ ہر صورت وچ
$R^+xz$، سو $R^+$ متعدی وی اے۔

«نالے» والی گل لئی، فرض کرو کہ $R$ باہم قابلِ موازنہ وی اے۔ سو ہر
$x \neq y$ لئی، یا $Rxy$ یا~$Ryx$، یعنی یا $\tuple{x, y} \in R$ یا
$\tuple{y, x} \in R$۔ کیوں جے $R \subseteq R^+$، ایہ گل $R^+$ لئی وی
درست رہندی اے، سو $R^+$ باہم قابلِ موازنہ وی اے۔
\end{proof}

\begin{prop}\ollabel{prop:partialtostrict}
جے $R$، سیٹ~$A$ اُتے جزوی ترتیب اے، تاں $R^- = R \setminus \Id{A}$ سخت
ترتیب اے۔ نالے، جے $R$ خطی ترتیب اے، تاں $R^-$ سخت خطی ترتیب اے۔
\end{prop}

\begin{proof}
ایہ ثبوت مشق دے طور تے چھڈّیا جاندا اے۔
\end{proof}

\begin{prob}
\olref[sfr][rel][ord]{prop:partialtostrict} دا ثبوت دیو۔
\end{prob}

اگلا سادہ نتیجہ ایہ ثابت کردا اے کہ سخت خطی ترتیباں، عنصراں نال تعیّن دے اصول
(Extensionality) ورگی اک خاصیت پوری کردیاں نیں:

\begin{prop}\ollabel{prop:extensionality-strictlinearorders}
جے $<$، سیٹ $A$ اُتے سخت خطی ترتیب اے، تاں:
\[
  (\forall a, b \in A)((\forall x \in A)(x < a \liff x < b) \lif a = b).
\]
\end{prop}

\begin{proof}
من لو کہ $(\forall x \in A)(x < a \liff x < b)$۔ جے $a < b$، تاں
$a < a$، جیہڑا ایس گل دے اُلٹ اے کہ $<$ اپنے آپ نال تعلق نہ ہون والا
اے؛ سو $a \nless b$۔ عین ایس طرح $b \nless a$۔ سو $a = b$، کیوں جے
$<$ باہم قابلِ موازنہ اے۔
\end{proof}

\end{document}
